\chapter{M\'odulo de Backups}

\section{Descripci\'on General}

El sistema de backups permite crear, listar, restaurar y eliminar copias
de seguridad de la base de datos PostgreSQL directamente desde la interfaz
web, utilizando las herramientas nativas \texttt{pg\_dump} y
\texttt{pg\_restore}.

\section{Utilidad de Backup (lib/backup.ts)}

El m\'odulo est\'a marcado como \texttt{'server-only'} para garantizar que
nunca se importe desde el cliente:

\begin{lstlisting}[style=typescript, caption=Backup utility - Creacion de backup]
import 'server-only';
import { exec } from 'child_process';
import { promisify } from 'util';
import fs from 'fs/promises';
import path from 'path';

const execAsync = promisify(exec);
const BACKUPS_DIR = path.join(process.cwd(), 'backups');
const DB_URL = process.env.DATABASE_URL!;

function parseDbUrl(url: string) {
  const u = new URL(url);
  return {
    host: u.hostname,
    port: u.port || '5432',
    user: decodeURIComponent(u.username),
    password: decodeURIComponent(u.password),
    database: u.pathname.replace('/', ''),
  };
}

export async function createBackup(): Promise<BackupInfo> {
  await ensureBackupsDir();

  const db = parseDbUrl(DB_URL);
  const timestamp = new Date().toISOString()
    .replace(/[:.]/g, '-').slice(0, 19);
  const filename = `vetrinaria-${timestamp}.dump`;
  const filepath = path.join(BACKUPS_DIR, filename);

  // Ejecuta pg_dump con formato personalizado (-Fc)
  await execAsync(
    `PGPASSWORD="${db.password}" pg_dump -h ${db.host} ` +
    `-p ${db.port} -U ${db.user} -d ${db.database} ` +
    `-Fc -f "${filepath}"`,
    { timeout: 300000 },
  );

  const stat = await fs.stat(filepath);

  return { filename, size: stat.size,
           createdAt: new Date().toISOString() };
}
\end{lstlisting}

\subsection{Restauraci\'on}

\begin{lstlisting}[style=typescript, caption=Restauracion de backup]
export async function restoreBackup(filename: string): Promise<void> {
  const db = parseDbUrl(DB_URL);
  const filepath = path.join(BACKUPS_DIR, filename);

  await fs.access(filepath);  // Verifica que existe

  // pg_restore con flag -c (clean: elimina objetos antes de restaurar)
  await execAsync(
    `PGPASSWORD="${db.password}" pg_restore -h ${db.host} ` +
    `-p ${db.port} -U ${db.user} -d ${db.database} ` +
    `-c -Fc "${filepath}"`,
    { timeout: 600000 },  // Timeout de 10 minutos
  );
}
\end{lstlisting}

\section{Server Actions}

\begin{lstlisting}[style=typescript, caption=Backups - Server actions]
// Trigger backup (acepta FormData para compatibilidad con formularios)
export async function triggerBackup(_formData: FormData): Promise<void> {
  await createBackup();
  revalidatePath('/settings/backups');
}

export async function getBackups(): Promise<BackupInfo[]> {
  return listBackups();
}

export async function removeBackup(filename: string) {
  await deleteBackup(filename);
  revalidatePath('/settings/backups');
}

export async function restoreFromBackup(filename: string) {
  await restoreBackup(filename);
  revalidatePath('/');
}
\end{lstlisting}

\section{Interfaz Web}

La p\'agina \texttt{/settings/backups} ofrece:

\begin{itemize}[leftmargin=*]
    \item Bot\'on \textbf{Crear Backup} que ejecuta pg\_dump y muestra
          el resultado.
    \item Lista de backups con nombre, tama\~no formateado (B, KB, MB),
          fecha de creaci\'on.
    \item Acciones por backup: \textbf{Restaurar} (con advertencia de
          sobreescritura en card \'ambar) y \textbf{Eliminar}.
    \item Panel de configuraci\'on de \textbf{Backups Autom\'aticos}
          con instrucciones para programar via cron.
    \item Estado vac\'{\i}o con mensaje y llamado a acci\'on.
\end{itemize}

\section{Backups Autom\'aticos (Cron)}

El script \texttt{scripts/backup.sh} puede ejecutarse via cron:

\begin{lstlisting}[style=bash, caption=Programar backup diario]
# Ejecutar a las 2:00 AM todos los dias
0 2 * * * /path/to/vetrinaria/scripts/backup.sh
\end{lstlisting}

\section{Seguridad}

\begin{itemize}[leftmargin=*]
    \item El m\'odulo \texttt{backup.ts} usa \texttt{'server-only'}.
    \item Las server actions verifican sesi\'on mediante middleware.
    \item La restauraci\'on muestra advertencia expl\'{\i}cita.
    \item Timeout de 5 minutos para creaci\'on y 10 para restauraci\'on.
\end{itemize}
